\chapter{Introduction générale}
\label{chap:introduction}

\section{Contexte}

Les simulations comportementales sont au cœur de nombreux domaines de
l'informatique~: jeux vidéo, films d'animation, recherche en biologie
artificielle, robotique en essaim, et applications grand public comme
les \emph{screensavers} animés. Elles cherchent à reproduire, par des
règles locales simples appliquées à un grand nombre d'agents, des
comportements globaux \emph{émergents} qui paraîtraient, sans cela,
extrêmement coûteux à scénariser explicitement.

Le cas particulier de la simulation d'un \emph{aquarium} concentre
plusieurs de ces problématiques~: déplacement fluide d'agents
hétérogènes, formation et dissolution de bancs, interactions
prédateur/proie, gestion d'une population qui naît et meurt, et tout
cela en temps réel avec un budget de calcul modeste. C'est un terrain
d'expérimentation idéal pour mettre en pratique le modèle
\emph{boids} de Reynolds~\cite{reynolds1987flocks}, qui reste, près de
quarante ans après sa publication, la référence pour ce type de
système.

Sur le plan technique, ce projet a également pour objectif de
développer une application graphique native sous Linux à l'aide de la
bibliothèque GTK4~\cite{gtk4docs}. Contrairement à un moteur de jeu
généraliste qui imposerait son propre cadre, GTK est un \emph{toolkit}
d'interface graphique de bureau~: l'utiliser pour piloter une
simulation animée oblige à exploiter astucieusement ses primitives
(positionnement absolu via \code{GtkFixed}, transformations CSS pour
la rotation, flux multimédia pour la vidéo de fond) et à respecter
son modèle d'événements basé sur la boucle principale GLib.

\section{Problématique}

Construire un aquarium virtuel crédible suppose de répondre à
plusieurs questions techniques imbriquées~:

\begin{enumerate}[leftmargin=*]
  \item \textbf{Comportement émergent.} Comment obtenir, à partir de
    règles locales (chaque poisson ne voit que ses voisins immédiats),
    un comportement global cohérent~: bancs qui se forment, se
    déplacent ensemble, évitent les obstacles, paniquent à l'approche
    d'un prédateur~?

  \item \textbf{Stabilité numérique.} Comment intégrer les forces qui
    s'exercent sur un poisson sans que la simulation diverge ou
    devienne saccadée~? La réponse passe par une boucle à pas fixe
    et un contrôle strict des accélérations et des vitesses.

  \item \textbf{Écosystème.} Comment représenter de façon généralisable
    les relations alimentaires entre espèces~? Comment les espèces
    naissent, meurent et se renouvellent~?

  \item \textbf{Adaptabilité.} Et surtout~: comment permettre à un
    utilisateur non programmeur de modifier l'écosystème sans
    recompiler le programme~? L'expérience montre qu'un système qui
    n'expose pas ses paramètres devient rapidement figé, donc
    inutilisable pour l'expérimentation.
\end{enumerate}

C'est sur ce dernier point que ce projet se distingue d'une
implémentation académique classique du modèle de Reynolds~: l'enjeu
n'est pas seulement de simuler, mais de fournir un système
\emph{ouvert} où la liste des espèces, leurs paramètres
comportementaux et les chaînes alimentaires sont décrits dans un
fichier texte modifiable à chaud.

\section{Objectifs}

Les objectifs assignés au projet sont les suivants~:

\begin{itemize}[leftmargin=*]
  \item Implémenter en C un moteur de simulation à pas fixe pour un
    écosystème de poissons, intégrant les trois forces de Reynolds
    (séparation, alignement, cohésion), une force de fuite face aux
    prédateurs, une force de chasse vers les proies, et une force
    d'évitement entre individus de la même espèce.

  \item Construire l'interface utilisateur avec GTK4~: fenêtre
    plein écran possible, vidéo de fond en boucle, contrôles latéraux
    (vitesse globale, pause, ajout et retrait précis de poissons,
    sélection des espèces du \emph{pool}, rechargement de la
    configuration).

  \item Concevoir un \textbf{système de configuration externe} en
    TOML~\cite{tomlspec} qui décrit l'intégralité du paramétrage~:
    toutes les espèces, toutes leurs propriétés, et les constantes
    globales du moteur. Cette configuration doit~:
    \begin{itemize}
      \item permettre l'ajout/suppression d'une espèce sans toucher
        au code C~;
      \item résoudre dynamiquement les références entre espèces
        (\code{prey} = liste d'identifiants)~;
      \item proposer des chemins de recherche standards et un
        mécanisme de \emph{fallback}~;
      \item être rechargeable à chaud depuis l'interface.
    \end{itemize}

  \item Garantir la \emph{robustesse} de l'application face à une
    configuration invalide~: erreurs de syntaxe TOML, identifiants
    manquants, références \code{prey} inconnues, valeurs hors plage.
    Le programme doit échouer proprement avec un message lisible
    plutôt que de produire un comportement indéfini.
\end{itemize}

\section{Plan du mémoire}

Ce rapport est organisé en sept chapitres~:

\begin{itemize}[leftmargin=*]
  \item Le \textbf{chapitre~\ref{chap:etat-de-lart}} présente l'état
    de l'art utile à la compréhension du projet~: le modèle de
    Reynolds, l'évolution vers des comportements de pilotage plus
    riches, les choix d'intégration numérique et un panorama
    comparatif des formats de configuration usuels.

  \item Le \textbf{chapitre~\ref{chap:analyse}} formalise les
    spécifications fonctionnelles et non fonctionnelles du système
    et présente les principaux cas d'utilisation.

  \item Le \textbf{chapitre~\ref{chap:conception}} décrit
    l'architecture modulaire retenue, le modèle de données, la
    séquence d'un \emph{tick} de simulation, le cycle de vie d'une
    entité et le graphe prédateur/proie.

  \item Le \textbf{chapitre~\ref{chap:implementation}} détaille
    l'implémentation~: environnement et chaîne de compilation,
    pipeline d'assets, algorithmes clés (forces de Reynolds, chasse,
    fuite, alimentation par la bouche, gestion des respawns), et
    enfin le système de configuration TOML.

  \item Le \textbf{chapitre~\ref{chap:tests}} présente les scénarios
    de test menés et leurs résultats.

  \item Le \textbf{chapitre~\ref{chap:conclusion}} dresse le bilan
    du projet et propose des pistes d'amélioration.

  \item Une bibliographie et plusieurs annexes (fichier de
    configuration complet, instructions de \emph{build},
    arborescence du projet, tableau récapitulatif des espèces)
    closent le document.
\end{itemize}
